\chapter{Instructions sur l'IA}

\section{Aperçu de la Politique}

Les outils de codage assistés par IA (ChatGPT, Claude, Copilot, Gemini,
Cursor, Codeium, et similaires) sont \textbf{autorisés} tout au long
de ce cursus. La recherche en sécurité est un domaine où la maîtrise
de l'IA devient essentielle --- apprendre à utiliser ces outils
efficacement fait partie de votre formation.

Cependant, l'objectif est la \textbf{compréhension}, pas l'automatisation.
Si vous ne pouvez pas expliquer chaque ligne de votre exploit sans
l'aide de l'IA, vous n'avez pas terminé l'exercice.

\begin{warning}
  Utiliser l'IA pour contourner l'objectif pédagogique d'un exercice
  --- par exemple, lui demander de rétro-ingéniérer un binaire pour
  ne pas avoir à apprendre le désassembleur, ou de générer une chaîne
  d'exploitation complète que vous ne comprenez pas --- va à l'encontre
  du but du cursus et sera traité comme une fraude académique.
\end{warning}

\section{Utilisations Autorisées}

\begin{itemize}
  \item \textbf{Explication de concepts}. Demander à l'IA de clarifier
        une classe de vulnérabilité, un sous-système du noyau, ou un
        protocole réseau.
  \item \textbf{Échafaudage de code}. Générer du code de départ pour
        un harnais, un fuzzer, ou un squelette de PoC --- puis le
        modifier et l'étendre vous-même.
  \item \textbf{Débogage}. Expliquer des messages d'erreur, la sortie
        d'un désassembleur, ou un comportement de crash inattendu.
  \item \textbf{Résumé de documentation}. Condenser une longue RFC ou
        un article technique en notes exploitables.
  \item \textbf{Traduction}. Traduire du contenu technique entre
        l'anglais et le français (les deux langues sont utilisées
        dans ce cursus).
  \item \textbf{Rédaction de rapports}. Demander de l'aide pour
        structurer un rapport de bug bounty --- mais remplir le
        contenu technique vous-même.
\end{itemize}

\section{Utilisations Interdites}

\begin{itemize}
  \item \textbf{Génération complète d'exploit}. Demander à l'IA
        d'écrire l'exploit complet d'un exercice, puis le soumettre
        sans le comprendre.
  \item \textbf{Contournement des contraintes de l'exercice}. Utiliser
        l'IA pour contourner une restriction que l'exercice impose
        explicitement (p.\ ex.\ «~n'utilisez pas de scanners
        automatisés~»).
  \item \textbf{Découverte automatisée de vulnérabilités}. Exécuter des
        outils de scan pilotés par IA qui remplacent l'analyse manuelle
        que vous êtes censé pratiquer.
  \item \textbf{Rédaction des grilles d'évaluation}. Vous n'êtes pas
        autorisé à utiliser l'IA pour générer des critères d'évaluation
        par les pairs ou des réponses.
  \item \textbf{Plagiat sans attribution}. Copier le résultat de l'IA
        sans le citer (voir Section~\ref{sec:ai-citation}).
\end{itemize}

\section{Responsabilité Personnelle}

Lorsque vous demandez à l'IA de générer du code d'exploit, vous
devenez responsable de chaque octet de ce résultat avant de
l'utiliser contre une cible.

\begin{enumerate}
  \item \textbf{Lisez le code généré} ligne par ligne.
  \item \textbf{Identifiez toutes les hypothèses} faites par l'IA
        (architecture, version du noyau, drapeaux du compilateur,
        topologie réseau).
  \item \textbf{Testez en isolation} avant d'exécuter contre une
        cible d'exercice.
  \item \textbf{Comprenez les modes d'échec}. Si l'exploit généré
        plante ou se bloque, vous devez diagnostiquer et corriger
        le problème vous-même.
\end{enumerate}

\hint{Si l'IA génère du code que vous ne comprenez pas, demandez-lui
de l'expliquer. Si vous ne comprenez toujours pas après
l'explication, le code n'est pas prêt à être utilisé. Passez à
une approche plus simple que vous maîtrisez entièrement.}

\section{Citation Obligatoire de l'IA}
\label{sec:ai-citation}

Chaque fichier contenant du code généré ou assisté par l'IA doit
inclure un commentaire de citation. Le format est~:

\begin{minted}{python}
# Generated with assistance from <tool> on <date>.
# Prompt: <one-line summary of the prompt>
# Verified by: <method used to verify the output>
\end{minted}

\subsection*{Exemples}

\textbf{Python}~:
\begin{minted}{python}
# Generated with assistance from Claude 3.5 Sonnet on 2025-03-14.
# Prompt: "Write a Python script that sends HTTP requests with
# modified price fields to /api/checkout/pay"
# Verified by: manual review of each request and response
\end{minted}

\textbf{C}~:
\begin{minted}{c}
/* Generated with assistance from GitHub Copilot on 2025-03-14.
 * Prompt: "Implement a stack-based buffer overflow with
 * ret2libc for x86_64 Linux"
 * Verified by: compiled with -fno-stack-protector and
 * tested against local binary with ASLR disabled */
\end{minted}

\textbf{Markdown (rapport)}~:
\begin{minted}{html}
<!--
Generated with assistance from ChatGPT-4 on 2025-03-14.
Prompt: "Structure a bug bounty report for IDOR with
sections: summary, steps, impact, mitigation"
Verified by: all technical content written manually,
AI only suggested structure
-->
\end{minted}

\subsection*{Interactions multiples avec l'IA}

Si vous avez utilisé plusieurs outils d'IA, ou itéré à travers
plusieurs invites, citez chaque interaction significative. Un
groupe d'invites connexes peut être résumé~:

\begin{minted}{html}
<!--
AI assistance used in this file:
- Claude 3.5 Sonnet (2025-03-14): harness skeleton generation
- ChatGPT-4 (2025-03-14): crash triage help
- Copilot (2025-03-15): inline autocomplete for loop logic
All generated code reviewed and modified manually.
-->
\end{minted}

\section{Évaluation par les Pairs et l'IA}

Lors de l'évaluation par les pairs, l'évaluateur peut~:

\begin{itemize}
  \item Vous demander d'\textbf{expliquer tout code généré par l'IA}
        ligne par ligne.
  \item Vous demander de \textbf{reproduire le raisonnement} sans
        l'aide de l'IA.
  \item Demander à voir votre \textbf{historique d'invites} pour un
        exercice donné.
  \item Déduire des points si la citation de l'IA est manquante ou
        si vous ne pouvez pas expliquer le code que l'IA a généré.
\end{itemize}

\begin{warning}
  Présenter le code généré par l'IA comme votre propre travail
  (défaut de citation) est du plagiat. Si l'évaluateur soupçonne
  du plagiat, il peut saisir l'équipe pédagogique pour examen.
\end{warning}

\section{Rappel Éthique}

\begin{itemize}
  \item Les modèles d'IA peuvent produire du code incorrect,
        dangereux, ou illégal. Vous êtes \textbf{responsable} de
        vérifier chaque résultat avant exécution.
  \item Les modèles d'IA entraînés sur la recherche publique en
        sécurité peuvent reproduire textuellement du code provenant
        de projets sous licence GPL. Vérifiez les licences avant
        d'intégrer ces éléments dans votre travail.
  \item Ne saisissez pas de code sensible ou propriétaire dans des
        outils d'IA hébergés dans le cloud. Utilisez des modèles
        locaux (Llama, Mistral, DeepSeek) pour les travaux impliquant
        des cibles non publiques.
  \item L'éthique de la recherche en sécurité que vous apprenez ici
        s'applique à l'utilisation de l'IA~: ne demandez pas à l'IA
        de générer des malwares, des rançongiciels, ou du code destiné
        à nuire à des systèmes en dehors du cadre de l'exercice.
\end{itemize}

\info{La maîtrise des outils d'IA est une compétence recherchée dans
la chasse aux bogues moderne. Utilisez-les sagement, citez-les
honnêtement, et conservez toujours votre propre compréhension de
chaque exploit que vous déployez.}
